\documentclass[11pt,a4paper]{article}
\usepackage[margin=2.25cm,headheight=22pt]{geometry}
\usepackage[T1]{fontenc}
\usepackage[utf8]{inputenc}
\usepackage[spanish,es-nodecimaldot]{babel}
\usepackage{lmodern}
\usepackage{microtype}
\usepackage{xcolor}
\usepackage{tabularx}
\usepackage{array}
\usepackage{enumitem}
\usepackage{fancyhdr}
\usepackage{hyperref}
\usepackage{titlesec}
\usepackage{tikz}
\usetikzlibrary{arrows.meta,positioning,babel}

\definecolor{azul}{HTML}{18324B}
\definecolor{turquesa}{HTML}{087E83}
\definecolor{gris}{HTML}{52616B}
\definecolor{fondo}{HTML}{EEF4F6}
\definecolor{linea}{HTML}{CFDCE1}

\hypersetup{
  colorlinks=true,
  linkcolor=turquesa,
  urlcolor=turquesa,
  pdftitle={Observando redes en mi entorno cotidiano},
  pdfauthor={Juan Skobalski}
}

\pagestyle{fancy}
\fancyhf{}
\fancyhead[L]{\small\color{gris} Ciberseguridad | Unidad 3 | Ejercicio 1}
\fancyhead[R]{\small\color{gris} Juan Skobalski}
\fancyfoot[C]{\small\color{gris}\thepage}
\renewcommand{\headrulewidth}{0.4pt}
\renewcommand{\headrule}{\hbox to\headwidth{\color{linea}\leaders\hrule height \headrulewidth\hfill}}

\titleformat{\section}[block]{\Large\bfseries\color{azul}}{}{0pt}{}
\titleformat{\subsection}[block]{\large\bfseries\color{turquesa}}{}{0pt}{}
\titlespacing*{\section}{0pt}{0pt}{8pt}
\titlespacing*{\subsection}{0pt}{14pt}{6pt}
\setlist[itemize]{leftmargin=1.35em,itemsep=4pt,topsep=4pt}
\setlist[enumerate]{leftmargin=1.55em,itemsep=5pt,topsep=4pt}
\setlength{\parskip}{6pt}
\setlength{\parindent}{0pt}
\renewcommand{\arraystretch}{1.28}

\newcommand{\infobox}[1]{%
  \noindent\colorbox{fondo}{\parbox{\dimexpr\textwidth-2\fboxsep}{#1}}}

\begin{document}

% Carátula
\thispagestyle{empty}
\vspace*{1.0cm}
{\color{turquesa}\large C O D E R H O U S E}\par
\vspace{0.4cm}
{\color{linea}\rule{\textwidth}{1.2pt}}\par
\vspace{2.3cm}
{\color{gris}\Large CIBERSEGURIDAD}\par
\vspace{0.45cm}
{\color{azul}\Huge\bfseries Observando redes en mi\\[5pt] entorno cotidiano}\par
\vspace{0.45cm}
{\color{turquesa}\LARGE Dispositivos, comunicación y paquetes de datos}\par
\vspace{1.0cm}
{\color{gris}\large Unidad 3 | Ejercicio 1}\par

\vfill
\noindent\colorbox{fondo}{%
  \begin{minipage}{\dimexpr\textwidth-2\fboxsep}
    \vspace{0.28cm}
    \noindent\makebox[3.5cm][l]{\textcolor{azul}{\textbf{Estudiante}}}Juan Skobalski\par
    \vspace{0.18cm}\color{linea}\hrule\color{black}\vspace{0.18cm}
    \noindent\makebox[3.5cm][l]{\textcolor{azul}{\textbf{Curso}}}Ciberseguridad\par
    \vspace{0.18cm}\color{linea}\hrule\color{black}\vspace{0.18cm}
    \noindent\makebox[3.5cm][l]{\textcolor{azul}{\textbf{Fecha}}}21 de octubre de 2026\par
    \vspace{0.28cm}
  \end{minipage}}
\vspace{0.75cm}
{\color{linea}\rule{\textwidth}{1.2pt}}\par
\vspace{0.25cm}
{\small\color{gris}Análisis de una red doméstica y del recorrido de los datos durante una actividad cotidiana.}

\newpage
\section*{Introducción}
\leavevmode\par

Una red es un conjunto de dispositivos conectados para intercambiar información. En el hogar, el router concentra las conexiones Wi-Fi, asigna direcciones IP privadas y dirige el tráfico dentro de la red local o hacia internet.

Este ejercicio identifica cinco dispositivos de uso cotidiano, analiza cómo pueden comunicarse y describe el recorrido de los datos al reproducir un video de YouTube desde un teléfono.

\subsection*{Objetivo y alcance}

El objetivo es distinguir entre comunicaciones que pueden permanecer dentro de la red local y aquellas que requieren servidores externos. El análisis representa un escenario doméstico general. En la lámpara y el lavarropas, el funcionamiento exacto depende de la aplicación y de la arquitectura elegida por cada fabricante.

\subsection*{Entorno observado}

Todos los equipos utilizan el punto de acceso Wi-Fi del router. Cada uno recibe una dirección privada, por ejemplo \texttt{192.168.1.10}, que lo identifica dentro de la LAN. Cuando el destino está en internet, el router aplica NAT y reenvía los paquetes al proveedor.

\begin{center}
\begin{tikzpicture}[
  node distance=8mm and 13mm,
  device/.style={draw=linea,fill=fondo,rounded corners=2pt,minimum width=3.3cm,minimum height=0.9cm,align=center,text=azul},
  router/.style={draw=turquesa,fill=turquesa!10,rounded corners=2pt,minimum width=3.6cm,minimum height=1.1cm,align=center,text=azul,font=\bfseries},
  link/.style={-{Latex[length=2mm]},draw=gris,thick}
]
  \node[router] (r) {Router doméstico\\Wi-Fi, LAN y NAT};
  \node[device,above left=of r] (mac) {MacBook Air M5};
  \node[device,above right=of r] (phone) {iPhone 17};
  \node[device,left=of r] (printer) {Impresora HP Inkjet};
  \node[device,right=of r] (lamp) {Lámpara Smart Home};
  \node[device,below=of r] (washer) {Lavarropas Midea};
  \draw[link] (mac) -- (r);
  \draw[link] (phone) -- (r);
  \draw[link] (printer) -- (r);
  \draw[link] (lamp) -- (r);
  \draw[link] (washer) -- (r);
\end{tikzpicture}
\end{center}

\infobox{\textcolor{azul}{\textbf{Idea central:}} estar conectado al mismo Wi-Fi permite formar parte de la misma red local, pero no garantiza que todas las funciones operen sin internet. Algunos fabricantes procesan las órdenes mediante servicios en la nube.}

\newpage
\section*{Inventario de dispositivos}

\begin{tabularx}{\textwidth}{>{\raggedright\arraybackslash\bfseries\color{azul}}p{3.3cm} >{\raggedright\arraybackslash}X >{\raggedright\arraybackslash}X}
  Dispositivo & Comunicación local & Uso de internet\\
  \hline
  MacBook Air M5 & Envía documentos a la impresora y puede intercambiar información con equipos autorizados. & Navegación, nube, actualizaciones y videollamadas.\\
  iPhone 17 & Descubre la impresora y puede controlar equipos inteligentes compatibles. & Mensajería, navegación, copias en la nube y contenido multimedia.\\
  Impresora HP Inkjet con Wi-Fi & Recibe trabajos desde la MacBook o el iPhone. & Internet es opcional para imprimir localmente; puede usarse para actualizaciones o servicios remotos.\\
  Lámpara Smart Home & Puede recibir órdenes locales si la aplicación y el producto lo admiten. & Puede depender de la nube para control remoto, automatizaciones o integración con una cuenta.\\
  Lavarropas Midea con Wi-Fi & Puede ser detectado o configurado desde la aplicación dentro del hogar. & Suele utilizar servicios externos para control remoto, avisos y actualizaciones.\\
\end{tabularx}

\subsection*{Comunicación dentro de la LAN}

Si la MacBook envía un documento a la impresora, ambos equipos intercambian paquetes mediante el router o su switch integrado. El tráfico puede permanecer dentro de la vivienda y la conexión del proveedor no es indispensable para transportar el trabajo de impresión.

\subsection*{Comunicación con internet}

Cuando el iPhone solicita un video o el lavarropas informa que terminó un ciclo mediante una plataforma remota, los paquetes salen hacia el proveedor. El router traduce la dirección privada mediante NAT y conserva la información necesaria para entregar la respuesta al dispositivo que inició la comunicación.

\subsection*{Dependencia de servicios externos}

Los dispositivos inteligentes no funcionan todos de la misma manera. Algunos permiten control local y otros envían cada orden a un servidor del fabricante. En el segundo caso, una interrupción de internet puede impedir el control desde la aplicación aunque el teléfono y el dispositivo continúen conectados al mismo Wi-Fi.

\newpage
\section*{Recorrido de los datos}

\subsection*{Acción elegida: reproducir YouTube en el iPhone}

El recorrido posible es el siguiente:

\begin{enumerate}
  \item El usuario abre YouTube en el iPhone y selecciona un video.
  \item El teléfono envía la solicitud por Wi-Fi al router doméstico.
  \item La aplicación consulta la dirección del servicio y establece una conexión cifrada.
  \item El router reenvía los paquetes al proveedor de internet.
  \item Los paquetes atraviesan redes intermedias hasta un servidor de YouTube o una red de distribución de contenido (CDN).
  \item El servidor responde enviando el video en múltiples partes y paquetes.
  \item La respuesta vuelve al proveedor, llega al router y se entrega al iPhone.
  \item El teléfono ordena los datos, reconstruye el contenido y comienza a reproducirlo mientras descarga las partes siguientes.
\end{enumerate}

\vspace{0.3cm}
\begin{center}
\begin{tikzpicture}[
  node distance=8mm,
  box/.style={draw=linea,fill=fondo,rounded corners=2pt,minimum width=7.3cm,minimum height=0.9cm,align=center,text=azul},
  link/.style={<->,>=Latex,draw=turquesa,thick}
]
  \node[box] (a) {iPhone 17};
  \node[box,below=of a] (b) {Router doméstico};
  \node[box,below=of b] (c) {Proveedor de internet};
  \node[box,below=of c] (d) {Redes intermedias};
  \node[box,below=of d] (e) {Servidor de YouTube o CDN};
  \draw[link] (a) -- node[right,font=\small]{Wi-Fi} (b);
  \draw[link] (b) -- (c);
  \draw[link] (c) -- (d);
  \draw[link] (d) -- (e);
\end{tikzpicture}
\end{center}

La comunicación es bidireccional: además de recibir video, el iPhone envía confirmaciones y nuevas solicitudes. El servidor puede ajustar la calidad según la velocidad disponible y continuar enviando las partes necesarias.

\newpage
\section*{Paquetes de datos}

La información no se transporta como un único bloque. Se divide en unidades pequeñas que contienen una parte de los datos y la información de control necesaria para reconocer su origen, destino y posición dentro de la comunicación.

\subsection*{Ventajas del sistema}

\begin{tabularx}{\textwidth}{>{\raggedright\arraybackslash\bfseries\color{azul}}p{4.0cm} >{\raggedright\arraybackslash}X}
  Uso eficiente & Los paquetes de diferentes usuarios y aplicaciones pueden compartir los mismos enlaces.\\
  Recuperación ante pérdidas & Si una parte no llega, se vuelve a enviar solamente la información afectada.\\
  Elección de rutas & Los paquetes pueden atravesar caminos distintos según la disponibilidad de las redes.\\
  Reconstrucción ordenada & La información de control permite ordenar las partes y reconstruir el contenido original.\\
  Escalabilidad & La infraestructura puede transportar páginas, mensajes, audio, video y archivos para una gran cantidad de equipos.\\
\end{tabularx}

\subsection*{Ejemplo aplicado}

Durante la reproducción de YouTube, la pérdida de un paquete no obliga a descargar nuevamente todo el video. Los mecanismos de transporte pueden recuperar la parte faltante, mientras otras partes continúan llegando y la aplicación utiliza su memoria intermedia para evitar interrupciones.

\subsection*{Conclusión}

La red doméstica combina comunicaciones locales y conexiones con servicios externos. La impresora puede funcionar principalmente dentro de la LAN, mientras que el iPhone, la computadora y los dispositivos inteligentes utilizan internet para muchas de sus funciones.

El router vincula ambos espacios: organiza el tráfico local y actúa como salida hacia otras redes. La división de la información en paquetes permite compartir esa infraestructura de manera eficiente, recuperar datos perdidos y entregar cada comunicación al destino correcto.

\infobox{\textcolor{azul}{\textbf{Síntesis:}} dispositivo, router, proveedor, redes intermedias y servidor forman el recorrido básico. La respuesta vuelve por la misma cadena lógica hasta que la aplicación reconstruye la información.}

\end{document}
